\documentclass[12pt]{article}

\usepackage{amsmath}
\usepackage{amssymb}
\usepackage{amsthm}
\usepackage{array}
\usepackage{geometry}
\usepackage{enumitem}
\usepackage{booktabs}

\geometry{
    margin=1in
}

\title{\textbf{APM1220: Applied Multivariate Data Analysis}\\
FA 1 -- Foundations: Matrix Algebra}
\author{Spike Lee-Roy V. Patayon}
\date{08/26/2026}

\begin{document}

\maketitle

\section*{Problem 2.2 -- Matrix Multiplication}

Given

\[
A=
\begin{bmatrix}
-1 & 3\\
4 & 2
\end{bmatrix},
\qquad
B=
\begin{bmatrix}
4 & -3\\
1 & -2\\
-2 & 0
\end{bmatrix},
\qquad
C=
\begin{bmatrix}
5\\
-4\\
2
\end{bmatrix}.
\]

The dimensions are

\[
A:2\times2,\qquad
B:3\times2,\qquad
C:3\times1.
\]

\subsection*{(a) $5A$}

Scalar multiplication does not change the dimensions of $A$.

\[
5A
=
5
\begin{bmatrix}
-1 & 3\\
4 & 2
\end{bmatrix}
\]

Multiplying every entry by $5$,

\[
\boxed{
5A=
\begin{bmatrix}
-5 & 15\\
20 & 10
\end{bmatrix}}
\]

\subsection*{(b) $BA$}

The dimensions are

\[
B_{3\times2}A_{2\times2}.
\]

The inner dimensions agree, so the multiplication is defined and the
result will be a $3\times2$ matrix.

\[
BA=
\begin{bmatrix}
4 & -3\\
1 & -2\\
-2 & 0
\end{bmatrix}
\begin{bmatrix}
-1 & 3\\
4 & 2
\end{bmatrix}.
\]

Calculate each entry:

\[
\begin{aligned}
(BA)_{11}
&=4(-1)+(-3)(4)
=-4-12
=-16,\\
(BA)_{12}
&=4(3)+(-3)(2)
=12-6
=6,\\
(BA)_{21}
&=1(-1)+(-2)(4)
=-1-8
=-9,\\
(BA)_{22}
&=1(3)+(-2)(2)
=3-4
=-1,\\
(BA)_{31}
&=(-2)(-1)+0(4)
=2,\\
(BA)_{32}
&=(-2)(3)+0(2)
=-6.
\end{aligned}
\]

Therefore,

\[
\boxed{
BA=
\begin{bmatrix}
-16 & 6\\
-9 & -1\\
2 & -6
\end{bmatrix}}
\]

\subsection*{(c) $A'B'$}

First calculate the transposes:

\[
A'=
\begin{bmatrix}
-1 & 4\\
3 & 2
\end{bmatrix}
\]

and

\[
B'=
\begin{bmatrix}
4 & 1 & -2\\
-3 & -2 & 0
\end{bmatrix}.
\]

The dimensions are

\[
A'_{2\times2}B'_{2\times3},
\]

so the result is $2\times3$.

Thus,

\[
A'B'
=
\begin{bmatrix}
-1 & 4\\
3 & 2
\end{bmatrix}
\begin{bmatrix}
4 & 1 & -2\\
-3 & -2 & 0
\end{bmatrix}.
\]

Calculate each entry:

\[
\begin{aligned}
(1,1)&=(-1)(4)+4(-3)=-16,\\
(1,2)&=(-1)(1)+4(-2)=-9,\\
(1,3)&=(-1)(-2)+4(0)=2,\\
(2,1)&=3(4)+2(-3)=6,\\
(2,2)&=3(1)+2(-2)=-1,\\
(2,3)&=3(-2)+2(0)=-6.
\end{aligned}
\]

Therefore,

\[
\boxed{
A'B'=
\begin{bmatrix}
-16 & -9 & 2\\
6 & -1 & -6
\end{bmatrix}}
\]

Notice that

\[
A'B'=(BA)'.
\]

\subsection*{(d) $C'B$}

First,

\[
C'=
\begin{bmatrix}
5 & -4 & 2
\end{bmatrix}.
\]

The dimensions are

\[
C'_{1\times3}B_{3\times2},
\]

so the result is $1\times2$.

\[
C'B=
\begin{bmatrix}
5 & -4 & 2
\end{bmatrix}
\begin{bmatrix}
4 & -3\\
1 & -2\\
-2 & 0
\end{bmatrix}.
\]

The first entry is

\[
5(4)+(-4)(1)+2(-2)
=20-4-4
=12.
\]

The second entry is

\[
5(-3)+(-4)(-2)+2(0)
=-15+8
=-7.
\]

Therefore,

\[
\boxed{
C'B=
\begin{bmatrix}
12 & -7
\end{bmatrix}}
\]

\subsection*{(e) Is $AB$ defined?}

The dimensions are

\[
A:2\times2
\]

and

\[
B:3\times2.
\]

For $AB$ to be defined, the inner dimensions must agree:

\[
(2\times\boxed{2})(\boxed{3}\times2).
\]

Since

\[
2\neq3,
\]

the multiplication is not defined.

Therefore,

\[
\boxed{\text{$AB$ is not defined.}}
\]

% --------------------------------------------------
\section*{Problem 2.4 -- Properties of Inverses}

Assume that $A^{-1}$ and $B^{-1}$ exist.

\subsection*{(a) Prove that $(A')^{-1}=(A^{-1})'$}

Since $A^{-1}$ exists,

\[
AA^{-1}=I.
\]

Taking the transpose of both sides gives

\[
(AA^{-1})'=I'.
\]

Using the property

\[
(XY)'=Y'X',
\]

we obtain

\[
(A^{-1})'A'=I.
\]

Also,

\[
A^{-1}A=I.
\]

Taking the transpose,

\[
(A^{-1}A)'=I',
\]

so

\[
A'(A^{-1})'=I.
\]

Therefore, $(A^{-1})'$ is both a left and right inverse of $A'$.
By uniqueness of the inverse,

\[
\boxed{
(A')^{-1}=(A^{-1})'
}
\]

\subsection*{(b) Prove that $(AB)^{-1}=B^{-1}A^{-1}$}

Consider

\[
(B^{-1}A^{-1})(AB).
\]

Using associativity,

\[
(B^{-1}A^{-1})(AB)
=
B^{-1}(A^{-1}A)B.
\]

Since

\[
A^{-1}A=I,
\]

we obtain

\[
B^{-1}IB
=
B^{-1}B
=
I.
\]

Now consider the multiplication in the other order:

\[
(AB)(B^{-1}A^{-1}).
\]

Again using associativity,

\[
(AB)(B^{-1}A^{-1})
=
A(BB^{-1})A^{-1}.
\]

Since

\[
BB^{-1}=I,
\]

we obtain

\[
AIA^{-1}
=
AA^{-1}
=
I.
\]

Thus $B^{-1}A^{-1}$ is the inverse of $AB$. Hence,

\[
\boxed{
(AB)^{-1}=B^{-1}A^{-1}
}
\]

% --------------------------------------------------
\section*{Problem 2.8 -- Eigenvalues, Normalized Eigenvectors,
and Spectral Decomposition}

Given

\[
A=
\begin{bmatrix}
1 & 2\\
2 & -2
\end{bmatrix}.
\]

The eigenvalues are found from

\[
|A-\lambda I|=0.
\]

Therefore,

\[
\begin{vmatrix}
1-\lambda & 2\\
2 & -2-\lambda
\end{vmatrix}
=0.
\]

Thus,

\[
(1-\lambda)(-2-\lambda)-4=0.
\]

Expanding,

\[
-2-\lambda+2\lambda+\lambda^2-4=0,
\]

so

\[
\lambda^2+\lambda-6=0.
\]

Factoring,

\[
(\lambda-2)(\lambda+3)=0.
\]

Therefore,

\[
\boxed{\lambda_1=2,\qquad\lambda_2=-3}.
\]

\subsection*{Eigenvector for $\lambda_1=2$}

Solve

\[
(A-2I)e_1=0.
\]

We have

\[
A-2I=
\begin{bmatrix}
-1 & 2\\
2 & -4
\end{bmatrix}.
\]

Therefore,

\[
-x+2y=0,
\]

so

\[
x=2y.
\]

Choose $y=1$. Then

\[
v_1=
\begin{bmatrix}
2\\
1
\end{bmatrix}.
\]

Its norm is

\[
\|v_1\|
=
\sqrt{2^2+1^2}
=
\sqrt{5}.
\]

Therefore, the normalized eigenvector is

\[
\boxed{
e_1=
\frac{1}{\sqrt{5}}
\begin{bmatrix}
2\\
1
\end{bmatrix}}
\]

\subsection*{Eigenvector for $\lambda_2=-3$}

Solve

\[
(A+3I)e_2=0.
\]

We have

\[
A+3I=
\begin{bmatrix}
4 & 2\\
2 & 1
\end{bmatrix}.
\]

Thus,

\[
4x+2y=0,
\]

which gives

\[
y=-2x.
\]

Choose $x=-1$, giving $y=2$. Therefore,

\[
v_2=
\begin{bmatrix}
-1\\
2
\end{bmatrix}.
\]

Its norm is

\[
\|v_2\|
=
\sqrt{(-1)^2+2^2}
=
\sqrt{5}.
\]

Hence,

\[
\boxed{
e_2=
\frac{1}{\sqrt{5}}
\begin{bmatrix}
-1\\
2
\end{bmatrix}}
\]

The eigenvectors are orthogonal because

\[
e_1'e_2
=
\frac{1}{5}
\begin{bmatrix}
2&1
\end{bmatrix}
\begin{bmatrix}
-1\\
2
\end{bmatrix}
=
\frac{-2+2}{5}
=0.
\]

\subsection*{Spectral Decomposition}

For a symmetric matrix,

\[
A=\lambda_1e_1e_1'+\lambda_2e_2e_2'.
\]

First,

\[
e_1e_1'
=
\frac{1}{5}
\begin{bmatrix}
2\\
1
\end{bmatrix}
\begin{bmatrix}
2&1
\end{bmatrix}
=
\frac{1}{5}
\begin{bmatrix}
4&2\\
2&1
\end{bmatrix}.
\]

Similarly,

\[
e_2e_2'
=
\frac{1}{5}
\begin{bmatrix}
-1\\
2
\end{bmatrix}
\begin{bmatrix}
-1&2
\end{bmatrix}
=
\frac{1}{5}
\begin{bmatrix}
1&-2\\
-2&4
\end{bmatrix}.
\]

Therefore,

\[
A
=
2
\left(
\frac{1}{5}
\begin{bmatrix}
4&2\\
2&1
\end{bmatrix}
\right)
-
3
\left(
\frac{1}{5}
\begin{bmatrix}
1&-2\\
-2&4
\end{bmatrix}
\right).
\]

Thus,

\[
A
=
\frac{1}{5}
\begin{bmatrix}
8&4\\
4&2
\end{bmatrix}
+
\frac{1}{5}
\begin{bmatrix}
-3&6\\
6&-12
\end{bmatrix}.
\]

Hence,

\[
A
=
\frac{1}{5}
\begin{bmatrix}
5&10\\
10&-10
\end{bmatrix}
=
\begin{bmatrix}
1&2\\
2&-2
\end{bmatrix}.
\]

Therefore,

\[
\boxed{
A=2e_1e_1'-3e_2e_2'
}
\]

where

\[
\boxed{
e_1=
\frac{1}{\sqrt5}
\begin{bmatrix}
2\\
1
\end{bmatrix},
\qquad
e_2=
\frac{1}{\sqrt5}
\begin{bmatrix}
-1\\
2
\end{bmatrix}.
}
\]

% --------------------------------------------------
\section*{Problem 2.11 -- Determinant of a Diagonal Matrix}

Let

\[
A=\{a_{ij}\}
\]

be a $p\times p$ diagonal matrix, so that

\[
a_{ij}=0,\qquad i\neq j.
\]

We want to prove that

\[
|A|=a_{11}a_{22}\cdots a_{pp}.
\]

Using cofactor expansion along the first row,

\[
|A|
=
a_{11}A_{11}
+a_{12}A_{12}
+\cdots
+a_{1p}A_{1p}.
\]

Since $A$ is diagonal,

\[
a_{12}=a_{13}=\cdots=a_{1p}=0.
\]

Therefore,

\[
|A|=a_{11}A_{11}.
\]

The cofactor $A_{11}$ is the determinant of the matrix obtained
by deleting the first row and first column:

\[
A_{11}
=
\begin{vmatrix}
a_{22}&0&\cdots&0\\
0&a_{33}&\cdots&0\\
\vdots&\vdots&\ddots&\vdots\\
0&0&\cdots&a_{pp}
\end{vmatrix}.
\]

This is again a diagonal matrix. Applying the same argument repeatedly,

\[
A_{11}=a_{22}a_{33}\cdots a_{pp}.
\]

Hence,

\[
|A|
=
a_{11}(a_{22}a_{33}\cdots a_{pp}).
\]

Therefore,

\[
\boxed{
|A|=a_{11}a_{22}\cdots a_{pp}
}
\]

or equivalently,

\[
\boxed{
|A|=\prod_{i=1}^{p}a_{ii}.
}
\]

% --------------------------------------------------
\section*{Problem 2.24 -- Covariance Matrix}

Given

\[
\Sigma=
\begin{bmatrix}
4&0&0\\
0&9&0\\
0&0&1
\end{bmatrix}.
\]

\subsection*{(a) Find $\Sigma^{-1}$}

Since $\Sigma$ is diagonal, its inverse is obtained by taking the
reciprocal of every nonzero diagonal entry:

\[
\boxed{
\Sigma^{-1}
=
\begin{bmatrix}
\frac14&0&0\\
0&\frac19&0\\
0&0&1
\end{bmatrix}}
\]

Verification:

\[
\Sigma\Sigma^{-1}
=
\begin{bmatrix}
4&0&0\\
0&9&0\\
0&0&1
\end{bmatrix}
\begin{bmatrix}
\frac14&0&0\\
0&\frac19&0\\
0&0&1
\end{bmatrix}
=
\begin{bmatrix}
1&0&0\\
0&1&0\\
0&0&1
\end{bmatrix}
=I.
\]

\subsection*{(b) Eigenvalues and eigenvectors of $\Sigma$}

Because $\Sigma$ is diagonal, its eigenvalues are its diagonal
entries:

\[
\boxed{
\lambda_1=4,\qquad
\lambda_2=9,\qquad
\lambda_3=1.
}
\]

The corresponding normalized eigenvectors are

\[
\boxed{
e_1=
\begin{bmatrix}
1\\0\\0
\end{bmatrix},
\qquad
e_2=
\begin{bmatrix}
0\\1\\0
\end{bmatrix},
\qquad
e_3=
\begin{bmatrix}
0\\0\\1
\end{bmatrix}.
}
\]

Thus,

\[
\boxed{
\begin{array}{c|c}
\text{Eigenvalue} & \text{Eigenvector}\\
\hline
4 & \begin{bmatrix}1\\0\\0\end{bmatrix}\\[6pt]
9 & \begin{bmatrix}0\\1\\0\end{bmatrix}\\[6pt]
1 & \begin{bmatrix}0\\0\\1\end{bmatrix}
\end{array}}
\]

\subsection*{(c) Eigenvalues and eigenvectors of $\Sigma^{-1}$}

We have

\[
\Sigma^{-1}
=
\begin{bmatrix}
\frac14&0&0\\
0&\frac19&0\\
0&0&1
\end{bmatrix}.
\]

Therefore,

\[
\boxed{
\lambda_1=\frac14,\qquad
\lambda_2=\frac19,\qquad
\lambda_3=1.
}
\]

The corresponding eigenvectors are unchanged:

\[
\boxed{
e_1=
\begin{bmatrix}
1\\0\\0
\end{bmatrix},
\qquad
e_2=
\begin{bmatrix}
0\\1\\0
\end{bmatrix},
\qquad
e_3=
\begin{bmatrix}
0\\0\\1
\end{bmatrix}.
}
\]

% --------------------------------------------------
\section*{Problem 2.30 -- Means, Covariances, and Linear Combinations}

Given

\[
X=
\begin{bmatrix}
X_1\\
X_2\\
X_3\\
X_4
\end{bmatrix},
\qquad
\mu_X=
\begin{bmatrix}
4\\
3\\
2\\
1
\end{bmatrix},
\]

and

\[
\Sigma_X=
\begin{bmatrix}
3&0&2&2\\
0&1&1&0\\
2&1&9&-2\\
2&0&-2&4
\end{bmatrix}.
\]

Partition $X$ as

\[
X=
\begin{bmatrix}
X^{(1)}\\
X^{(2)}
\end{bmatrix},
\]

where

\[
X^{(1)}
=
\begin{bmatrix}
X_1\\
X_2
\end{bmatrix},
\qquad
X^{(2)}
=
\begin{bmatrix}
X_3\\
X_4
\end{bmatrix}.
\]

Let

\[
A=
\begin{bmatrix}
1&2
\end{bmatrix},
\qquad
B=
\begin{bmatrix}
1&-2\\
2&-1
\end{bmatrix}.
\]

Partitioning the mean vector,

\[
\mu_{X^{(1)}}=
\begin{bmatrix}
4\\
3
\end{bmatrix},
\qquad
\mu_{X^{(2)}}=
\begin{bmatrix}
2\\
1
\end{bmatrix}.
\]

The covariance blocks are

\[
\Sigma_{11}
=
\operatorname{Cov}(X^{(1)})
=
\begin{bmatrix}
3&0\\
0&1
\end{bmatrix},
\]

\[
\Sigma_{22}
=
\operatorname{Cov}(X^{(2)})
=
\begin{bmatrix}
9&-2\\
-2&4
\end{bmatrix},
\]

and

\[
\Sigma_{12}
=
\operatorname{Cov}(X^{(1)},X^{(2)})
=
\begin{bmatrix}
2&2\\
1&0
\end{bmatrix}.
\]

\subsection*{(a) $E(X^{(1)})$}

\[
\boxed{
E(X^{(1)})
=
\begin{bmatrix}
4\\
3
\end{bmatrix}}
\]

\subsection*{(b) $E(AX^{(1)})$}

Using

\[
E(AX^{(1)})
=
AE(X^{(1)}),
\]

we obtain

\[
E(AX^{(1)})
=
\begin{bmatrix}
1&2
\end{bmatrix}
\begin{bmatrix}
4\\
3
\end{bmatrix}.
\]

Therefore,

\[
E(AX^{(1)})
=
4+6
=
10.
\]

Hence,

\[
\boxed{
E(AX^{(1)})=10
}
\]

\subsection*{(c) $\operatorname{Cov}(X^{(1)})$}

From the upper-left $2\times2$ block of $\Sigma_X$,

\[
\boxed{
\operatorname{Cov}(X^{(1)})
=
\begin{bmatrix}
3&0\\
0&1
\end{bmatrix}}
\]

\subsection*{(d) $\operatorname{Cov}(AX^{(1)})$}

Use

\[
\operatorname{Cov}(AX^{(1)})
=
A\operatorname{Cov}(X^{(1)})A'.
\]

Thus,

\[
\operatorname{Cov}(AX^{(1)})
=
\begin{bmatrix}
1&2
\end{bmatrix}
\begin{bmatrix}
3&0\\
0&1
\end{bmatrix}
\begin{bmatrix}
1\\
2
\end{bmatrix}.
\]

First,

\[
\begin{bmatrix}
1&2
\end{bmatrix}
\begin{bmatrix}
3&0\\
0&1
\end{bmatrix}
=
\begin{bmatrix}
3&2
\end{bmatrix}.
\]

Therefore,

\[
\operatorname{Cov}(AX^{(1)})
=
\begin{bmatrix}
3&2
\end{bmatrix}
\begin{bmatrix}
1\\
2
\end{bmatrix}
=
3+4
=
7.
\]

Hence,

\[
\boxed{
\operatorname{Cov}(AX^{(1)})=[7]
}
\]

\subsection*{(e) $E(X^{(2)})$}

\[
\boxed{
E(X^{(2)})
=
\begin{bmatrix}
2\\
1
\end{bmatrix}}
\]

\subsection*{(f) $E(BX^{(2)})$}

Using

\[
E(BX^{(2)})
=
BE(X^{(2)}),
\]

we obtain

\[
E(BX^{(2)})
=
\begin{bmatrix}
1&-2\\
2&-1
\end{bmatrix}
\begin{bmatrix}
2\\
1
\end{bmatrix}.
\]

The first component is

\[
1(2)-2(1)=0.
\]

The second component is

\[
2(2)-1(1)=3.
\]

Therefore,

\[
\boxed{
E(BX^{(2)})
=
\begin{bmatrix}
0\\
3
\end{bmatrix}}
\]

\subsection*{(g) $\operatorname{Cov}(X^{(2)})$}

From the lower-right $2\times2$ block of $\Sigma_X$,

\[
\boxed{
\operatorname{Cov}(X^{(2)})
=
\begin{bmatrix}
9&-2\\
-2&4
\end{bmatrix}}
\]

\subsection*{(h) $\operatorname{Cov}(BX^{(2)})$}

Use

\[
\operatorname{Cov}(BX^{(2)})
=
B\Sigma_{22}B'.
\]

We have

\[
B'=
\begin{bmatrix}
1&2\\
-2&-1
\end{bmatrix}.
\]

First,

\[
B\Sigma_{22}
=
\begin{bmatrix}
1&-2\\
2&-1
\end{bmatrix}
\begin{bmatrix}
9&-2\\
-2&4
\end{bmatrix}.
\]

Multiplying,

\[
B\Sigma_{22}
=
\begin{bmatrix}
13&-10\\
20&-8
\end{bmatrix}.
\]

Therefore,

\[
\operatorname{Cov}(BX^{(2)})
=
\begin{bmatrix}
13&-10\\
20&-8
\end{bmatrix}
\begin{bmatrix}
1&2\\
-2&-1
\end{bmatrix}.
\]

Thus,

\[
\operatorname{Cov}(BX^{(2)})
=
\begin{bmatrix}
33&36\\
36&48
\end{bmatrix}.
\]

Hence,

\[
\boxed{
\operatorname{Cov}(BX^{(2)})
=
\begin{bmatrix}
33&36\\
36&48
\end{bmatrix}}
\]

\subsection*{(i) $\operatorname{Cov}(X^{(1)},X^{(2)})$}

This is the upper-right block of $\Sigma_X$:

\[
\boxed{
\operatorname{Cov}(X^{(1)},X^{(2)})
=
\begin{bmatrix}
2&2\\
1&0
\end{bmatrix}}
\]

\subsection*{(j) $\operatorname{Cov}(AX^{(1)},BX^{(2)})$}

For linear transformations,

\[
\operatorname{Cov}(AX^{(1)},BX^{(2)})
=
A\operatorname{Cov}(X^{(1)},X^{(2)})B'.
\]

The dimensions are

\[
(1\times2)(2\times2)(2\times2)
\rightarrow1\times2.
\]

Thus,

\[
\operatorname{Cov}(AX^{(1)},BX^{(2)})
=
\begin{bmatrix}
1&2
\end{bmatrix}
\begin{bmatrix}
2&2\\
1&0
\end{bmatrix}
\begin{bmatrix}
1&2\\
-2&-1
\end{bmatrix}.
\]

First,

\[
\begin{bmatrix}
1&2
\end{bmatrix}
\begin{bmatrix}
2&2\\
1&0
\end{bmatrix}
=
\begin{bmatrix}
4&2
\end{bmatrix}.
\]

Then,

\[
\begin{bmatrix}
4&2
\end{bmatrix}
\begin{bmatrix}
1&2\\
-2&-1
\end{bmatrix}
=
\begin{bmatrix}
0&6
\end{bmatrix}.
\]

Therefore,

\[
\boxed{
\operatorname{Cov}(AX^{(1)},BX^{(2)})
=
\begin{bmatrix}
0&6
\end{bmatrix}}
\]

% --------------------------------------------------
\section*{Problem 2.41 -- Mean and Covariance of Linear Combinations}

Given

\[
\mu_X=
\begin{bmatrix}
3\\
2\\
-2\\
0
\end{bmatrix}
\]

and

\[
\Sigma_X=
\begin{bmatrix}
3&0&0&0\\
0&3&0&0\\
0&0&3&0\\
0&0&0&3
\end{bmatrix}
=
3I_4.
\]

Also,

\[
A=
\begin{bmatrix}
1&-1&0&0\\
1&1&-2&0\\
1&1&1&-3
\end{bmatrix}.
\]

The dimensions are

\[
A:3\times4,
\qquad
X:4\times1,
\]

so

\[
AX:3\times1.
\]

The three linear combinations are

\[
L_1=X_1-X_2,
\]

\[
L_2=X_1+X_2-2X_3,
\]

and

\[
L_3=X_1+X_2+X_3-3X_4.
\]

\subsection*{(a) Find $E(AX)$}

Using

\[
E(AX)=AE(X)=A\mu_X,
\]

we obtain

\[
E(AX)
=
\begin{bmatrix}
1&-1&0&0\\
1&1&-2&0\\
1&1&1&-3
\end{bmatrix}
\begin{bmatrix}
3\\
2\\
-2\\
0
\end{bmatrix}.
\]

For the first component,

\[
1(3)-1(2)=1.
\]

For the second component,

\[
1(3)+1(2)-2(-2)
=
3+2+4
=
9.
\]

For the third component,

\[
1(3)+1(2)+1(-2)-3(0)
=
3+2-2
=
3.
\]

Therefore,

\[
\boxed{
E(AX)
=
\begin{bmatrix}
1\\
9\\
3
\end{bmatrix}}
\]

\subsection*{(b) Find $\operatorname{Cov}(AX)$}

For a linear transformation,

\[
\operatorname{Cov}(AX)
=
A\Sigma_XA'.
\]

Since

\[
\Sigma_X=3I_4,
\]

we have

\[
\operatorname{Cov}(AX)
=
3AA'.
\]

Calculate $AA'$.

The first row of $A$ is

\[
r_1=
\begin{bmatrix}
1&-1&0&0
\end{bmatrix}.
\]

The second row is

\[
r_2=
\begin{bmatrix}
1&1&-2&0
\end{bmatrix}.
\]

The third row is

\[
r_3=
\begin{bmatrix}
1&1&1&-3
\end{bmatrix}.
\]

The diagonal entries are

\[
r_1r_1'
=
1^2+(-1)^2
=
2,
\]

\[
r_2r_2'
=
1^2+1^2+(-2)^2
=
6,
\]

and

\[
r_3r_3'
=
1^2+1^2+1^2+(-3)^2
=
12.
\]

For the off-diagonal entries,

\[
r_1r_2'
=
1(1)+(-1)(1)+0(-2)+0(0)
=
0,
\]

\[
r_1r_3'
=
1(1)+(-1)(1)+0(1)+0(-3)
=
0,
\]

and

\[
r_2r_3'
=
1(1)+1(1)+(-2)(1)+0(-3)
=
0.
\]

Therefore,

\[
AA'
=
\begin{bmatrix}
2&0&0\\
0&6&0\\
0&0&12
\end{bmatrix}.
\]

Hence,

\[
\operatorname{Cov}(AX)
=
3
\begin{bmatrix}
2&0&0\\
0&6&0\\
0&0&12
\end{bmatrix}.
\]

Thus,

\[
\boxed{
\operatorname{Cov}(AX)
=
\begin{bmatrix}
6&0&0\\
0&18&0\\
0&0&36
\end{bmatrix}}
\]

Therefore, the variances are

\[
\boxed{
\operatorname{Var}(L_1)=6,
\qquad
\operatorname{Var}(L_2)=18,
\qquad
\operatorname{Var}(L_3)=36.
}
\]

The covariances are

\[
\boxed{
\operatorname{Cov}(L_1,L_2)=0,
}
\]

\[
\boxed{
\operatorname{Cov}(L_1,L_3)=0,
}
\]

and

\[
\boxed{
\operatorname{Cov}(L_2,L_3)=0.
}
\]

\subsection*{(c) Which pairs of linear combinations have zero covariances?}

From

\[
\operatorname{Cov}(AX)
=
\begin{bmatrix}
6&0&0\\
0&18&0\\
0&0&36
\end{bmatrix},
\]

all off-diagonal entries are zero.

Therefore,

\[
\boxed{
\operatorname{Cov}(L_1,L_2)=0
}
\]

\[
\boxed{
\operatorname{Cov}(L_1,L_3)=0
}
\]

and

\[
\boxed{
\operatorname{Cov}(L_2,L_3)=0.
}
\]

Hence, every pair of the three linear combinations has zero covariance.

% --------------------------------------------------
\section*{Final Answers Summary}

\subsection*{Problem 2.2}

\[
\boxed{
5A=
\begin{bmatrix}
-5&15\\
20&10
\end{bmatrix}}
\]

\[
\boxed{
BA=
\begin{bmatrix}
-16&6\\
-9&-1\\
2&-6
\end{bmatrix}}
\]

\[
\boxed{
A'B'=
\begin{bmatrix}
-16&-9&2\\
6&-1&-6
\end{bmatrix}}
\]

\[
\boxed{
C'B=
\begin{bmatrix}
12&-7
\end{bmatrix}}
\]

\[
\boxed{\text{$AB$ is not defined.}}
\]

\subsection*{Problem 2.4}

\[
\boxed{
(A')^{-1}=(A^{-1})'
}
\]

\[
\boxed{
(AB)^{-1}=B^{-1}A^{-1}
}
\]

\subsection*{Problem 2.8}

\[
\boxed{\lambda_1=2}
\]

\[
\boxed{
e_1=
\frac{1}{\sqrt5}
\begin{bmatrix}
2\\1
\end{bmatrix}}
\]

\[
\boxed{\lambda_2=-3}
\]

\[
\boxed{
e_2=
\frac{1}{\sqrt5}
\begin{bmatrix}
-1\\2
\end{bmatrix}}
\]

\[
\boxed{
A=2e_1e_1'-3e_2e_2'
}
\]

\subsection*{Problem 2.11}

\[
\boxed{
|A|
=
a_{11}a_{22}\cdots a_{pp}
=
\prod_{i=1}^{p}a_{ii}
}
\]

\subsection*{Problem 2.24}

\[
\boxed{
\Sigma^{-1}
=
\begin{bmatrix}
\frac14&0&0\\
0&\frac19&0\\
0&0&1
\end{bmatrix}}
\]

Eigenvalues of $\Sigma$:

\[
\boxed{4,\;9,\;1}
\]

Eigenvalues of $\Sigma^{-1}$:

\[
\boxed{\frac14,\;\frac19,\;1}
\]

The corresponding eigenvectors are the standard basis vectors.

\subsection*{Problem 2.30}

\[
\boxed{
E(X^{(1)})
=
\begin{bmatrix}
4\\3
\end{bmatrix}}
\]

\[
\boxed{
E(AX^{(1)})=10
}
\]

\[
\boxed{
\operatorname{Cov}(X^{(1)})
=
\begin{bmatrix}
3&0\\
0&1
\end{bmatrix}}
\]

\[
\boxed{
\operatorname{Cov}(AX^{(1)})=[7]
}
\]

\[
\boxed{
E(X^{(2)})
=
\begin{bmatrix}
2\\1
\end{bmatrix}}
\]

\[
\boxed{
E(BX^{(2)})
=
\begin{bmatrix}
0\\3
\end{bmatrix}}
\]

\[
\boxed{
\operatorname{Cov}(X^{(2)})
=
\begin{bmatrix}
9&-2\\
-2&4
\end{bmatrix}}
\]

\[
\boxed{
\operatorname{Cov}(BX^{(2)})
=
\begin{bmatrix}
33&36\\
36&48
\end{bmatrix}}
\]

\[
\boxed{
\operatorname{Cov}(X^{(1)},X^{(2)})
=
\begin{bmatrix}
2&2\\
1&0
\end{bmatrix}}
\]

\[
\boxed{
\operatorname{Cov}(AX^{(1)},BX^{(2)})
=
\begin{bmatrix}
0&6
\end{bmatrix}}
\]

\subsection*{Problem 2.41}

\[
\boxed{
E(AX)
=
\begin{bmatrix}
1\\9\\3
\end{bmatrix}}
\]

\[
\boxed{
\operatorname{Cov}(AX)
=
\begin{bmatrix}
6&0&0\\
0&18&0\\
0&0&36
\end{bmatrix}}
\]

All three pairs have zero covariance:

\[
\boxed{
\operatorname{Cov}(L_1,L_2)=0,
\qquad
\operatorname{Cov}(L_1,L_3)=0,
\qquad
\operatorname{Cov}(L_2,L_3)=0.
}
\]

\end{document}